\documentclass[11pt]{article}

\usepackage[margin=1in]{geometry}
\usepackage{amsmath,amssymb,amsthm}
\usepackage{booktabs}
\usepackage{microtype}
\usepackage{natbib}
\usepackage[hidelinks]{hyperref}

\newtheorem{proposition}{Proposition}
\newtheorem{corollary}{Corollary}
\newtheorem{lemma}{Lemma}
\newtheorem{example}{Example}
\newcommand{\clip}{\operatorname{clip}}
\newcommand{\DC}{\Delta c}

\title{History-Based Price Discrimination with Switching Costs:\\
Equilibrium Scope and a Welfare Correction}
\author{Anonymous draft}
\date{\today}

\begin{document}
\maketitle

\begin{abstract}
We revisit the Hotelling model of history-based price discrimination in
\citet{gehrig2011}. Two corrections matter for its equilibrium and welfare
comparison. First, once cohort demands are globally clipped, the displayed
interior price formulas need not be global Nash equilibria: a firm charging one
price across customer cohorts may profitably raise price and abandon one cohort.
We derive explicit parameter domains under which the displayed
history-based-pricing and uniform-pricing profiles are Nash equilibria and show
that the uniform-pricing validity domain is contained in the
history-based-pricing domain. Second, direct integration of consumer surplus
yields a factor of three missing from the complete author draft's reported
consumer-surplus difference. On the common validity domain,
uniform pricing raises consumer surplus by
$3\theta(1-\theta)\sigma^2/(16\tau)$, while the incumbent's profit loss is only
$\theta(1-\theta)\sigma^2/(8\tau)$ and the entrant's profit is unchanged.
Hence total welfare is higher under the displayed uniform-pricing equilibrium
by $\theta(1-\theta)\sigma^2/(16\tau)$. The welfare correction has a
real-resource interpretation: history-based pricing improves spatial matching,
but the resulting increase in switching-cost losses is larger. We do not claim
game-wide uniqueness or characterize all equilibria outside the displayed
profiles' validity domains.
\end{abstract}

\noindent\textbf{Keywords:} history-based price discrimination; switching costs;
Hotelling competition; global equilibrium; welfare.

\section{Introduction}

\citet{gehrig2011} study history-based price discrimination (HBP) in a market
with an incumbent, an entrant, inherited customers, new customers, and
switching costs. Their model produces a sharp comparison between HBP and
uniform pricing, including the conclusions that entry is invariant across the
two pricing regimes and that the incumbent's gain from HBP exceeds the
associated consumer loss.

This note re-evaluates two logically distinct parts of that comparison. The
first is equilibrium scope. The displayed prices solve the smooth first-order
conditions when both customer cohorts are interior. But a common-price firm
faces a piecewise objective once demands are clipped to the unit interval. A
finite price increase can move the firm to a different active set---for
example, serving only the higher-value cohort---and can dominate the smooth
candidate. We characterize parameter domains on which the displayed HBP and
uniform-pricing profiles are global Nash equilibria.

The second correction concerns welfare accounting. Direct integration of the
primitive utilities yields a consumer-surplus difference three times the
expression reported in the complete author draft. Once that factor is restored,
the total-welfare comparison reverses on the domain where both displayed price
profiles are valid equilibria. The correction is economically substantive, not
a transfer-accounting artifact: HBP reduces transportation mismatch but
increases switching-cost losses by a larger amount.

Our contribution is deliberately source-specific. We do not develop a new
general theory of segmented Bertrand competition, and we do not claim that
uniform pricing is generally welfare superior to history-based or personalized
pricing. Related work studies broader or different environments, including
asymmetric no-discrimination constraints \citep{bouckaert2013}, models in which
both firms can discriminate \citep{gehrig2012,gnutzmann2014}, dynamic
history-based pricing \citep{shy2016}, and welfare under richer information
structures \citep{chen2026}. The present results concern the model and displayed
price profiles in \citet{gehrig2011}.

\paragraph{Source version.}
Our equation-level reconstruction uses the authors' complete 14 December 2009
draft \citep{gehrig2009draft}. We separately verified the bibliographic metadata
and published abstract for the 2011 article. The published abstract itself
retains the welfare headline that the incumbent's gain from HBP exceeds the
associated consumer loss. We therefore distinguish two source statements:
the missing factor in the displayed consumer-surplus expression is attributed
to the complete author draft, while the disputed welfare headline is also
present in the published article record. We do not claim line-by-line
inspection of the Version of Record equations.

\section{Model and displayed price profiles}

Firm $A$, the incumbent, is located at $0$ and firm $B$, the entrant, at $1$ on
a unit Hotelling line. A mass $1-\theta$ of consumers is inherited and a mass
$\theta$ is new, with $0<\theta<1$. Transportation cost is $\tau>0$. An old
consumer who buys from $B$ incurs switching cost $\sigma\geq 0$. Marginal costs
are $c_A$ and $c_B$, and we write
\[
  \DC \equiv c_A-c_B.
\]
As in the source model, we maintain full market coverage; clipping below refers
to allocation between the two firms rather than market participation.

The complete author draft does not explicitly state whether price strategies
belong to $\mathbb R$, $\mathbb R_+$, or another set. We state the equilibrium
results for the real-price game. If prices are instead restricted to be
nonnegative, the same global-deviation conditions apply whenever the displayed
candidate prices and the relevant deviations are feasible; in particular this
covers the standard nonnegative-cost, positive-markup regular case.

Under HBP, $A$ charges $p_A$ to old consumers and $q_A$ to new consumers,
whereas $B$ charges a common price $p_B$. Old-consumer utilities are
\[
 u_A^o(x)=\beta-p_A-\tau x,\qquad
 u_B^o(x)=\beta-p_B-\tau(1-x)-\sigma,
\]
and new-consumer utilities are
\[
 u_A^n(x)=\beta-q_A-\tau x,\qquad
 u_B^n(x)=\beta-p_B-\tau(1-x).
\]
Thus the actual incumbent shares are
\[
 x_o=\clip\!\left(\frac12+\frac{\sigma+p_B-p_A}{2\tau},0,1\right),
 \qquad
 x_n=\clip\!\left(\frac12+\frac{p_B-q_A}{2\tau},0,1\right).
\]
The corresponding operating profits are
\begin{align}
 \pi_A&=(1-\theta)(p_A-c_A)x_o
       +\theta(q_A-c_A)x_n,\\
 \pi_B&=(p_B-c_B)\bigl[(1-\theta)(1-x_o)+\theta(1-x_n)\bigr].
\end{align}
Under uniform pricing, $A$ also charges one price, so $p_A=q_A$.

The displayed smooth HBP candidate is
\begin{align}
 p_A^d &=
 \tau+\frac{(2+\theta)\sigma+4c_A+2c_B}{6},\\
 q_A^d &=
 \tau+\frac{4c_A+2c_B-(1-\theta)\sigma}{6},\\
 p_B^d &=
 \tau+\frac{2c_B+c_A-(1-\theta)\sigma}{3},
\end{align}
with shares
\begin{align}
 x_o^d&=\frac12+\frac{(2+\theta)\sigma-2\DC}{12\tau},\\
 x_n^d&=\frac12-\frac{(1-\theta)\sigma+2\DC}{12\tau}.
\end{align}
The displayed uniform-pricing candidate is
\begin{align}
 p_A^u&=\tau+\frac{2c_A+c_B+(1-\theta)\sigma}{3},\\
 p_B^u&=\tau+\frac{2c_B+c_A-(1-\theta)\sigma}{3},
\end{align}
with
\begin{align}
 x_o^u&=\frac12+\frac{(2\theta+1)\sigma-\DC}{6\tau},\\
 x_n^u&=\frac12-\frac{\DC+2(1-\theta)\sigma}{6\tau}.
\end{align}
Both regimes imply the same aggregate incumbent share on the displayed branch:
\[
  m_A^d=m_A^u
  =\frac12+\frac{(1-\theta)\sigma-\DC}{6\tau}.
\]

These formulas solve the interior first-order conditions. The next section
shows why that fact alone is not a global-equilibrium certificate.

\section{Global validity of the displayed equilibria}

The nontrivial issue is common-price coupling. Once one cohort drops out, the
profit function changes branch, and a second local maximum can appear.

\begin{proposition}[Global validity of the displayed profiles]
\label{prop:validity}
Let $0<\theta<1$, $\tau>0$, and $\sigma\geq0$.

\begin{enumerate}
\item The displayed HBP profile is a global Nash equilibrium on
\[
\mathcal D_d:\quad
-3\tau+\frac{4-\theta+3\sqrt{\theta}}{4}\sigma
\leq \DC
\leq
3\tau-\frac{1-\theta}{2}\sigma.
\]

\item The displayed uniform-pricing profile is a global Nash equilibrium if and
only if
\[
\mathcal D_u:\quad
-3\tau+\frac{2+\theta+3\sqrt{\theta}}{2}\sigma
\leq \DC
\leq
3\tau-\frac{1-\theta+3\sqrt{1-\theta}}{2}\sigma
\]
for the displayed-candidate construction.
\end{enumerate}
\end{proposition}

The lower HBP bound prevents the entrant from abandoning old consumers and
moving to its higher-price new-only branch. In the uniform-pricing game the
lower bound similarly prevents a new-only deviation by $B$, while the upper
bound prevents $A$ from raising price and serving only old consumers. The
proof, given in Appendix~\ref{app:global}, explicitly optimizes every clipped
branch.

\begin{corollary}[Common comparison domain]
\label{cor:containment}
For $0<\theta<1$ and $\sigma\geq0$,
\[
  \mathcal D_u\subseteq \mathcal D_d.
\]
Hence the common validity domain for comparing the two displayed equilibria is
\[
  \mathcal D_{\mathrm{compare}}=\mathcal D_u.
\]
\end{corollary}

Indeed, the uniform lower bound exceeds the HBP lower bound by
\[
  \frac34\bigl(\theta+\sqrt{\theta}\bigr)\sigma,
\]
while the uniform upper bound lies below the HBP upper bound by
\[
  \frac32\sqrt{1-\theta}\,\sigma.
\]

\begin{example}[Interior shares do not imply global Nash]
\label{ex:hbp}
Take
\[
 \theta=\frac12,\quad \tau=1,\quad \sigma=\frac{23}{10},
 \quad c_A=c_B=0.
\]
The displayed HBP candidate is
\[
 p_A=\frac{47}{24},\qquad
 q_A=\frac{97}{120},\qquad
 p_B=\frac{37}{60},
\]
with strictly interior incumbent shares
\[
 x_o=\frac{47}{48},\qquad x_n=\frac{97}{240}.
\]
The entrant's candidate profit is $1369/7200$. By instead charging
$p_B'=217/240$, the entrant abandons old consumers and obtains
$47089/230400$, an exact gain of
\[
 \frac{3281}{230400}>0.
\]
Thus strict cohort interiority is insufficient for global equilibrium.
\end{example}

At weak boundaries of Proposition~\ref{prop:validity}, tied best responses can
occur. We therefore make no game-wide uniqueness claim and no
selection-invariant welfare claim for alternative boundary equilibria.

\section{Consumer surplus and welfare correction}

Let $CS^r$ and $\pi_i^r$ denote consumer surplus and firm $i

\begin{proposition}[Corrected welfare comparison]
\label{prop:welfare}
On $\mathcal D_{\mathrm{compare}}$,
\begin{align}
 CS^u-CS^d
 &=\frac{3\theta(1-\theta)\sigma^2}{16\tau},\\
 \pi_A^d-\pi_A^u
 &=\frac{\theta(1-\theta)\sigma^2}{8\tau},\\
 \pi_B^d-\pi_B^u&=0,
\end{align}
and therefore
\begin{equation}
 W^d-W^u
 =-\frac{\theta(1-\theta)\sigma^2}{16\tau}.
\end{equation}
For $0<\theta<1$, $\sigma>0$, and $\tau>0$, total welfare is strictly higher
at the displayed uniform-pricing equilibrium.
\end{proposition}

The complete author draft reports the same consumer-surplus expression without
the factor $3$. Restoring that factor changes the total-welfare calculation:
the incumbent's profit gain from HBP equals only two thirds of the consumer
loss, while entrant profit is unchanged.

Table~\ref{tab:welfare} summarizes the decomposition in units of
$\theta(1-\theta)\sigma^2/\tau$.

\begin{table}[ht]
\centering
\caption{HBP minus uniform pricing: distributional and resource decomposition}
\label{tab:welfare}
\begin{tabular}{lr}
\toprule
Component & Coefficient in HBP $-$ uniform \\
\midrule
Consumer surplus & $-3/16$ \\
Incumbent profit & $+2/16$ \\
Entrant profit & $0$ \\
\midrule
Total welfare & $-1/16$ \\
\midrule
Transportation cost & $-3/16$ \\
Switching-cost loss & $+4/16$ \\
Total real resource cost & $+1/16$ \\
\bottomrule
\end{tabular}
\end{table}

\section{Why the welfare ranking reverses}

The allocation changes in a particularly transparent way:
\begin{align}
 x_o^d-x_o^u&=-\frac{\theta\sigma}{4\tau},\\
 x_n^d-x_n^u&=\frac{(1-\theta)\sigma}{4\tau}.
\end{align}
Thus HBP sends more old consumers to entrant $B$ and more new consumers to
incumbent $A$, while leaving the aggregate incumbent share unchanged. Prices
move in parallel:
\begin{align}
 p_A^d-p_A^u&=\frac{\theta\sigma}{2},\\
 q_A^d-p_A^u&=-\frac{(1-\theta)\sigma}{2},\\
 p_B^d-p_B^u&=0.
\end{align}

For a cohort with incumbent cutoff $x$, aggregate Hotelling transportation cost
is
\[
 \tau\left(x^2-x+\frac12\right).
\]
Substitution yields
\[
 TC^d-TC^u
 =-\frac{3\theta(1-\theta)\sigma^2}{16\tau}.
\]
HBP therefore improves spatial matching. But switching costs are borne only by
old consumers who buy from $B$, and
\[
 SC^d-SC^u
 =\frac{\theta(1-\theta)\sigma^2}{4\tau}
 =\frac{4\theta(1-\theta)\sigma^2}{16\tau}.
\]
The increase in switching losses exceeds the transportation-cost saving,
leaving
\[
 (TC^d+SC^d)-(TC^u+SC^u)
 =\frac{\theta(1-\theta)\sigma^2}{16\tau}.
\]
This exactly equals the negative of the welfare difference in
Proposition~\ref{prop:welfare}.

This resource interpretation assumes that $\sigma$ is a real switching
disutility or resource cost---for example time, hassle, learning, or
compatibility loss. A purely monetary termination fee may instead be partly a
transfer and need not map one-for-one into this welfare accounting.

\section{Scope and conclusion}

The note establishes two source-specific corrections. First, the displayed
interior pricing formulas require explicit global-validity restrictions because
common pricing across cohorts creates profitable active-set deviations outside
those domains. Second, on the common certified domain, correcting consumer
surplus reverses the displayed-equilibrium welfare ranking.

Several claims are intentionally outside our scope. We do not characterize the
complete equilibrium correspondence outside $\mathcal D_d$ or $\mathcal D_u$;
we do not claim uniqueness of the clipped price game; and we do not infer an
optimal regulatory policy or endogenous regime choice. Likewise, the welfare
result is not a general theorem that uniform pricing dominates personalized or
history-based pricing in other models. We also do not challenge the
entry-invariance identity merely because the welfare calculation is corrected:
on the common displayed branch the entrant-profit difference remains zero, so
any entry statement must still be evaluated conditional on the validity of the
compared equilibria and the source's entry-cost comparison.

The surviving result is narrower and sharper: within the
\citet{gehrig2011} environment, the displayed comparison is valid only on an
explicit parameter domain, and on that domain the corrected welfare accounting
favors uniform pricing because the additional switching loss induced by HBP
exceeds its spatial-matching gain.

\appendix

\section{Global best-response calculation}
\label{app:global}

Normalize prices by $\tau$, recenter by a firm's marginal cost, and let $z$ be
its normalized net markup. For a firm charging one common price to two cohorts,
write normalized demand as
\[
 D(z)
 =w\clip\!\left(\frac{h_H-z}{2},0,1\right)
 +(1-w)\clip\!\left(\frac{h_L-z}{2},0,1\right),
 \qquad h_H\geq h_L.
\]
Let
\[
 \bar h=wh_H+(1-w)h_L.
\]
When both positive-mass segments are interior, the smooth FOC gives
$z^*=\bar h/2$ and profit $\bar h^2/8$.

\begin{lemma}[Two-segment common-price globality]
Suppose the displayed smooth candidate has both segments interior and positive
profit. Its only potentially superior high-price local maximum is the
high-segment-only vertex, whose peak profit is $wh_H^2/8$. Consequently the
smooth candidate is a global best response if and only if
\[
 \bar h^2\geq wh_H^2,
\]
or, on the positive regular branch,
\[
 \bar h\geq\sqrt{w}\,h_H.
\]
\end{lemma}

\begin{proof}
The clipped demand is continuous and piecewise affine, so
$\varphi(z)=zD(z)$ is continuous and piecewise quadratic. Because both
segments are interior at $z^*=\bar h/2$ and the candidate profit is positive,
$z^*>0$. On the both-interior branch,
\[
 \varphi'(z)=\frac{\bar h-2z}{2},
\]
which is positive below $z^*$ and negative above it.

Consider first deviations below $z^*$. Every breakpoint encountered when
moving upward toward the both-interior branch is a transition of a segment from
full demand to interior demand. At such a breakpoint $D$ is continuous and its
slope becomes weakly more negative. Hence, for nonnegative $z$, the one-sided
derivative $D(z)+zD'(z)$ weakly falls when crossing the breakpoint from left to
right. Since it is positive immediately below $z^*$, working backward through
all lower-price branches shows that it remains positive there. Prices with
negative net markup give nonpositive profit and cannot dominate the positive
candidate.

Above $z^*$, the only possible source of a new local maximum is a transition
at which the lower-threshold segment's demand falls to zero. Thereafter the
remaining high-threshold segment yields
\[
 \varphi_H(z)=\frac{w}{2}z(h_H-z),
\]
whose unconstrained vertex is $h_H/2$ with value $wh_H^2/8$. If that vertex
does not lie on the high-only branch, the feasible branch maximum occurs at its
boundary and the comparison below is automatically satisfied under the
both-interior feasibility conditions. Thus the only potentially superior
high-price peak is bounded by $wh_H^2/8$.

The both-interior candidate yields $\bar h^2/8$. Comparing these two peaks
gives $\bar h^2\geq wh_H^2$. On the positive regular branch this is equivalent
to $\bar h\geq\sqrt w\,h_H$.
\end{proof}

For HBP, the incumbent's two prices solve separate one-segment problems. The
entrant is the common-price firm, with the new cohort as the high-threshold
segment of weight $\theta$. Substitution in the lemma yields the lower bound in
$\mathcal D_d$; the binding incumbent segment condition yields its upper bound.

Under uniform pricing, both firms face common-price coupling. For $B$, the
new cohort is the high-threshold segment of weight $\theta$, yielding the lower
bound in $\mathcal D_u$. For $A$, the old cohort is the high-threshold segment
of weight $1-\theta$, yielding the upper bound.

\section{Additional exact deviations}

The two sides of the uniform-pricing validity domain also admit exact
counterexamples with strictly interior displayed shares.

For an entrant-side failure, take
\[
 \theta=\frac12,\quad \tau=\sigma=1,\quad c_B=0,\quad c_A=-\frac45.
\]
The displayed prices are $p_A=19/30$ and $p_B=17/30$, with
$x_o=29/30$ and $x_n=7/15$. A new-only entrant deviation to $49/60$ yields an
exact profit gain of
\[
 \frac{89}{14400}>0.
\]

For an incumbent-side failure, take
\[
 \theta=\frac12,\quad \tau=\sigma=1,\quad c_B=0,\quad c_A=\frac95.
\]
The displayed prices are $p_A=71/30$ and $p_B=43/30$, with
$x_o=8/15$ and $x_n=1/30$. An old-only incumbent deviation to $157/60$ again
yields an exact profit gain of
\[
 \frac{89}{14400}>0.
\]

The first example can be translated upward by a common additive shift in both
marginal costs and all prices without changing demand or markups, so the
negative value of $c_A$ is not economically essential to the deviation.

\section{Verification and reproducibility}

The proof-critical algebra and exact counterexamples are independently checked
in the accompanying repository. A clean-room verifier reconstructs clipped
payoffs directly from primitives and enumerates all piecewise-quadratic
unilateral-deviation branches. A separate Lean~4 development verifies the
threshold reductions, domain containment, welfare identities, welfare sign, and
exact rational counterexamples.

The formalization intentionally does not encode the full game or the
exhaustiveness of the economic case partition; those objects are certified by
the independent global-deviation audit rather than inferred from the proof
assistant. The manuscript makes no claim stronger than that layered
certification.

\bibliographystyle{apalike}
\bibliography{references}

\end{document}
s operating profit
under regime $r\in\{d,u\}$. For any displayed allocation
$(x_o,x_n)$ and prices $(p_A,q_A,p_B)$, aggregate consumer surplus is
\begin{align}
CS={}&(1-\theta)\left[
 \int_0^{x_o}(\beta-p_A-\tau x)\,dx
 +\int_{x_o}^1(\beta-p_B-\tau(1-x)-\sigma)\,dx
 \right] \\
&+\theta\left[
 \int_0^{x_n}(\beta-q_A-\tau x)\,dx
 +\int_{x_n}^1(\beta-p_B-\tau(1-x))\,dx
 \right].
\label{eq:csprimitive}
\end{align}
Under uniform pricing set $q_A=p_A$. Substituting the two displayed price and
allocation profiles into \eqref{eq:csprimitive} gives the following result.

\begin{proposition}[Corrected welfare comparison]
\label{prop:welfare}
On $\mathcal D_{\mathrm{compare}}$,
\begin{align}
 CS^u-CS^d
 &=\frac{3\theta(1-\theta)\sigma^2}{16\tau},\\
 \pi_A^d-\pi_A^u
 &=\frac{\theta(1-\theta)\sigma^2}{8\tau},\\
 \pi_B^d-\pi_B^u&=0,
\end{align}
and therefore
\begin{equation}
 W^d-W^u
 =-\frac{\theta(1-\theta)\sigma^2}{16\tau}.
\end{equation}
For $0<\theta<1$, $\sigma>0$, and $\tau>0$, total welfare is strictly higher
at the displayed uniform-pricing equilibrium.
\end{proposition}

The complete author draft reports the same consumer-surplus expression without
the factor $3$. Restoring that factor changes the total-welfare calculation:
the incumbent's profit gain from HBP equals only two thirds of the consumer
loss, while entrant profit is unchanged.

Table~\ref{tab:welfare} summarizes the decomposition in units of
$\theta(1-\theta)\sigma^2/\tau$.

\begin{table}[ht]
\centering
\caption{HBP minus uniform pricing: distributional and resource decomposition}
\label{tab:welfare}
\begin{tabular}{lrr}
\toprule
Component & HBP $-$ uniform & Coefficient \\
\midrule
Consumer surplus & $-3/16$ & $-3/16$ \\
Incumbent profit & $+2/16$ & $+2/16$ \\
Entrant profit & $0$ & $0$ \\
\midrule
Total welfare & $-1/16$ & $-1/16$ \\
\midrule
Transportation cost & $-3/16$ & $-3/16$ \\
Switching-cost loss & $+4/16$ & $+4/16$ \\
Total real resource cost & $+1/16$ & $+1/16$ \\
\bottomrule
\end{tabular}
\end{table}

\section{Why the welfare ranking reverses}

The allocation changes in a particularly transparent way:
\begin{align}
 x_o^d-x_o^u&=-\frac{\theta\sigma}{4\tau},\\
 x_n^d-x_n^u&=\frac{(1-\theta)\sigma}{4\tau}.
\end{align}
Thus HBP sends more old consumers to entrant $B$ and more new consumers to
incumbent $A$, while leaving the aggregate incumbent share unchanged. Prices
move in parallel:
\begin{align}
 p_A^d-p_A^u&=\frac{\theta\sigma}{2},\\
 q_A^d-p_A^u&=-\frac{(1-\theta)\sigma}{2},\\
 p_B^d-p_B^u&=0.
\end{align}

For a cohort with incumbent cutoff $x$, aggregate Hotelling transportation cost
is
\[
 \tau\left(x^2-x+\frac12\right).
\]
Substitution yields
\[
 TC^d-TC^u
 =-\frac{3\theta(1-\theta)\sigma^2}{16\tau}.
\]
HBP therefore improves spatial matching. But switching costs are borne only by
old consumers who buy from $B$, and
\[
 SC^d-SC^u
 =\frac{\theta(1-\theta)\sigma^2}{4\tau}
 =\frac{4\theta(1-\theta)\sigma^2}{16\tau}.
\]
The increase in switching losses exceeds the transportation-cost saving,
leaving
\[
 (TC^d+SC^d)-(TC^u+SC^u)
 =\frac{\theta(1-\theta)\sigma^2}{16\tau}.
\]
This exactly equals the negative of the welfare difference in
Proposition~\ref{prop:welfare}.

This resource interpretation assumes that $\sigma$ is a real switching
disutility or resource cost---for example time, hassle, learning, or
compatibility loss. A purely monetary termination fee may instead be partly a
transfer and need not map one-for-one into this welfare accounting.

\section{Scope and conclusion}

The note establishes two source-specific corrections. First, the displayed
interior pricing formulas require explicit global-validity restrictions because
common pricing across cohorts creates profitable active-set deviations outside
those domains. Second, on the common certified domain, correcting consumer
surplus reverses the displayed-equilibrium welfare ranking.

Several claims are intentionally outside our scope. We do not characterize the
complete equilibrium correspondence outside $\mathcal D_d$ or $\mathcal D_u$;
we do not claim uniqueness of the clipped price game; and we do not infer an
optimal regulatory policy or endogenous regime choice. Likewise, the welfare
result is not a general theorem that uniform pricing dominates personalized or
history-based pricing in other models.

The surviving result is narrower and sharper: within the
\citet{gehrig2011} environment, the displayed comparison is valid only on an
explicit parameter domain, and on that domain the corrected welfare accounting
favors uniform pricing because the additional switching loss induced by HBP
exceeds its spatial-matching gain.

\appendix

\section{Global best-response calculation}
\label{app:global}

Normalize prices by $\tau$, recenter by a firm's marginal cost, and let $z$ be
its normalized net markup. For a firm charging one common price to two cohorts,
write normalized demand as
\[
 D(z)
 =w\clip\!\left(\frac{h_H-z}{2},0,1\right)
 +(1-w)\clip\!\left(\frac{h_L-z}{2},0,1\right),
 \qquad h_H\geq h_L.
\]
Let
\[
 \bar h=wh_H+(1-w)h_L.
\]
When both positive-mass segments are interior, the smooth FOC gives
$z^*=\bar h/2$ and profit $\bar h^2/8$.

\begin{lemma}[Two-segment common-price globality]
Suppose the displayed smooth candidate has both segments interior and positive
profit. Its only potentially superior high-price local maximum is the
high-segment-only vertex, whose peak profit is $wh_H^2/8$. Consequently the
smooth candidate is a global best response if and only if
\[
 \bar h^2\geq wh_H^2,
\]
or, on the positive regular branch,
\[
 \bar h\geq\sqrt{w}\,h_H.
\]
\end{lemma}

\begin{proof}
The clipped demand is continuous and piecewise affine, so $zD(z)$ is continuous
and piecewise quadratic. Moving upward from the both-interior branch, the only
new quadratic branch after the lower-threshold cohort exits is
$wz(h_H-z)/2$, with vertex $h_H/2$ and value $wh_H^2/8$. The lower-price
branches cannot dominate the both-interior vertex: on the branch adjacent from
below, the one-sided derivative at the breakpoint is at least the derivative
on the both-interior branch and is positive before $z^*$; the full-demand
branch is linear in the markup. Negative-markup deviations yield nonpositive
profit. Comparing the two relevant peaks gives the stated condition.
\end{proof}

For HBP, the incumbent's two prices solve separate one-segment problems. The
entrant is the common-price firm, with the new cohort as the high-threshold
segment of weight $\theta$. Substitution in the lemma yields the lower bound in
$\mathcal D_d$; the binding incumbent segment condition yields its upper bound.

Under uniform pricing, both firms face common-price coupling. For $B$, the
new cohort is the high-threshold segment of weight $\theta$, yielding the lower
bound in $\mathcal D_u$. For $A$, the old cohort is the high-threshold segment
of weight $1-\theta$, yielding the upper bound.

\section{Additional exact deviations}

The two sides of the uniform-pricing validity domain also admit exact
counterexamples with strictly interior displayed shares.

For an entrant-side failure, take
\[
 \theta=\frac12,\quad \tau=\sigma=1,\quad c_B=0,\quad c_A=-\frac45.
\]
The displayed prices are $p_A=19/30$ and $p_B=17/30$, with
$x_o=29/30$ and $x_n=7/15$. A new-only entrant deviation to $49/60$ yields an
exact profit gain of
\[
 \frac{89}{14400}>0.
\]

For an incumbent-side failure, take
\[
 \theta=\frac12,\quad \tau=\sigma=1,\quad c_B=0,\quad c_A=\frac95.
\]
The displayed prices are $p_A=71/30$ and $p_B=43/30$, with
$x_o=8/15$ and $x_n=1/30$. An old-only incumbent deviation to $157/60$ again
yields an exact profit gain of
\[
 \frac{89}{14400}>0.
\]

The first example can be translated upward by a common additive shift in both
marginal costs and all prices without changing demand or markups, so the
negative value of $c_A$ is not economically essential to the deviation.

\section{Verification and reproducibility}

The proof-critical algebra and exact counterexamples are independently checked
in the accompanying repository. A clean-room verifier reconstructs clipped
payoffs directly from primitives and enumerates all piecewise-quadratic
unilateral-deviation branches. A separate Lean~4 development verifies the
threshold reductions, domain containment, welfare identities, welfare sign, and
exact rational counterexamples.

The formalization intentionally does not encode the full game or the
exhaustiveness of the economic case partition; those objects are certified by
the independent global-deviation audit rather than inferred from the proof
assistant. The manuscript makes no claim stronger than that layered
certification.

\bibliographystyle{apalike}
\bibliography{references}

\end{document}
